\section{The Second Bill for a Shareable \code{node}}
\label{sec:ch15-the-second-bill-for-a-shareable-node}

\index{node@\code{node(id:)}!what the field has now cost twice|(}%
Chapter~\ref{ch:composition} shipped a type interceptor of thirty-three lines
so that \code{node} and \code{nodes} would carry \code{@shareable} and two
subgraphs would compose. I argued the trade at the time: thirty-three lines to
avoid promising four fields in order to keep a promise about two.

\index{node@\code{node(id:)}!the runtime bill, from chapter 9}%
Chapter~\ref{ch:second-third-subgraph} presented the first bill. With two node
types in two services, \code{node(id:)} answers correctly only when the query
names the owning type in a fragment, because the fragment is the only thing in
the request the planner can route on. I said there that the graph ships as
measured and that chapter~\ref{ch:satisfiability} would pick it up.

\index{node@\code{node(id:)}!the composition bill}%
This is the second bill, and it is larger, because the first one is at least
visible to whoever sends the query.
Section~\ref{sec:ch15-the-composition-that-passes} is the same field, the same
directive and the same two subgraphs, costing a check that was supposed to be
protecting the whole graph rather than one root field.

\index{node@\code{node(id:)}!one mechanism, two costs}%
It is one mechanism twice. \code{@shareable} says any subgraph declaring the
field can resolve it, and neither of them can. At run time the router believes
it and picks one. At composition the walk believes it and accepts either.

\subsection{One of the Four Rejections, Measured}

\index{node@\code{node(id:)}!declaring it in one subgraph only}%
Chapter~\ref{ch:second-third-subgraph} listed four ways out of the runtime
problem and rejected all four. The second was to declare \code{node} in one
subgraph only, so that nothing collides and no \code{@shareable} is needed, and
I rejected it on the ground that composition would pass and the failure would
be identical. That was reasoning rather than a measurement, and the
measurement disagrees with half of it.

\index{node@\code{node(id:)}!the check comes back}%
Take \code{node} and \code{nodes} out of the Speakers document, leave
everything else alone, and the healthy graph composes exactly as it does today.
Break the Ratings key again, with all four services in the graph, and the
composer catches it, naming
\code{query \{ node \{ ... on Session \} \}}: the route
section~\ref{sec:ch15-the-route-it-named-is-one-of-six} got by taking the whole
Speakers service out, now reached by taking two fields out of it. So that
option is not neutral. It costs nothing at composition time on a working graph,
and it buys back the check.

\index{node@\code{node(id:)}!what it would cost at run time is not known}%
What it costs at run time I have not measured. With one subgraph declaring the
field, a \code{node(id:)} for a speaker has to be planned through the Sessions
subgraph's copy, and whether that answers, errors or reaches the wrong service
is a question I did not put to the router. I am not going to tell you it is the
better arrangement on the strength of half a comparison.

\subsection{Nobody Has Answered This}

\index{node id!there is no published guidance, and it is not an oversight}%
Chapter~\ref{ch:federation-model} recorded that this book federates on the
global node id as my judgment, because Apollo publishes no guidance on using
one as a key. That is still true and it is worth saying more precisely now.
The question has been open on Apollo's own federation repository since October
2021, under the title \enquote{Node query support in federated apollo server},
and it was still open, with eleven comments on it, when I wrote this
sentence~\autocite{apollonodeissue}.

\index{node id!five years is a position}%
Five years of an issue staying open is not the same as nobody having noticed.
Relay's global object identification and federation's entity resolution are two
answers to the question \emph{how do I refetch a thing by its identifier}, and
a graph that adopts both has to decide which of them routes. This book decided,
in chapter~\ref{ch:schema-design}, before there was a seam to route across.

\index{node id!what I would tell you to do}%
What I would do differently, having now paid for it twice: keep global object
identification, because chapter~\ref{ch:schema-design}'s refetch story is worth
having, and try declaring \code{node} in exactly one subgraph before sharing it
across every subgraph that generates it. That is a judgment and I am labeling
it as one, on the evidence in this section rather than on anything published.
It buys back the check that section~\ref{sec:ch15-the-composition-that-passes}
loses, at a runtime cost I have not measured and you would have to. Which is
the honest version of a recommendation whose other half is still open on
somebody else's issue tracker.
\index{node@\code{node(id:)}!what the field has now cost twice|)}%
